\documentclass[11pt,a4paper]{article}

% --- Encoding and fonts ---
\usepackage[T1]{fontenc}
\usepackage[utf8]{inputenc}
\usepackage{lmodern}
\usepackage{microtype}

% --- Math ---
\usepackage{amsmath}
\usepackage{amssymb}
\usepackage{amsthm}
\usepackage{mathtools}

% --- Layout and tables ---
\usepackage[a4paper,margin=1in]{geometry}
\usepackage{array}
\usepackage{booktabs}
\usepackage{enumitem}
\usepackage{xcolor}

% Ragged paragraph columns and a breakable monospace for narrow table cells
\newcolumntype{L}[1]{>{\raggedright\arraybackslash}p{#1}}
\newcommand{\bcode}[1]{\texttt{\small #1}}

% --- Links (load last) ---
\usepackage[colorlinks=true,linkcolor=blue!55!black,citecolor=blue!55!black,
            urlcolor=blue!55!black,breaklinks=true]{hyperref}

% --- Theorem environments ---
\theoremstyle{plain}
\newtheorem{lemma}{Lemma}

% --- Convenience macros ---
\newcommand{\CHO}{\mathbb{C}\otimes\mathbb{H}\otimes\mathbb{O}}
\newcommand{\Alg}{\mathcal{A}}
\newcommand{\JO}{\mathfrak{J}_3(\mathbb{O})}
\newcommand{\OP}{\mathbb{OP}^2}
\newcommand{\ffour}{\mathfrak{f}_4}
\newcommand{\Spin}{\mathrm{Spin}}
\newcommand{\Der}{\mathrm{Der}}
\newcommand{\Cl}{\mathrm{Cl}}
\newcommand{\code}[1]{\texttt{#1}}
\newcommand{\epso}{\varepsilon_0}

\hypersetup{
  pdftitle  = {What C x H x O and J3(O) Fix and Do Not Fix About One Standard-Model Generation and the Family Count},
  pdfauthor = {Richard Astbury},
  pdfkeywords = {division algebras; octonions; exceptional Jordan algebra; Standard Model;
                 three generations; KO-dimension; reproducible computation}
}

\title{\bfseries What $\CHO$ and $\JO$ Fix---and Do Not Fix---About\\
One Standard-Model Generation and the Family Count\\[4pt]
\large A conservative, falsifiable delimitation}

\author{Richard Astbury\thanks{Contact: \href{mailto:richard.astbury@gmail.com}{richard.astbury@gmail.com}.
Developed with the assistance of AI tools. The author is solely
responsible for all claims, code, and errors.}}

\date{Version 1.0 \quad (content frozen 7 June 2026; typeset June 2026)}

\begin{document}
\maketitle

\begin{abstract}
\noindent
We state and defend a deliberately narrow claim about $\Alg = \CHO$, the tensor
product of the division algebras $\mathbb{C}$, $\mathbb{H}$, and $\mathbb{O}$
(64 real dimensions). We note at once that $\Alg$ is \emph{not} itself a division
algebra---tensor products of division algebras carry zero divisors (already
$\mathbb{C}\otimes\mathbb{O}$, the complex octonions, does)---so we never use that
term for it. We do \emph{not} claim a
theory of the Standard Model's continuous parameters, nor a theory of everything.
We claim only that the \emph{discrete internal structure} of one fermion
generation---its electric charges and colour multiplicities, its weak-isospin and
hypercharge assignments, its gauge-anomaly cancellation, and its chirality---is
determined, once a complex structure is fixed, by $\Alg$ and its automorphisms,
with no per-observable continuous input. The qualifier ``once a complex structure
is fixed'' is load-bearing: that choice is \emph{adopted} (from Furey), not
derived, and is the single most consequential of the discrete choices below, so
this part of the claim is in substance Furey's result, which we inherit rather
than reclaim. The family count \emph{three} enters
separately, as the \emph{rank} of the exceptional Jordan algebra $\JO$, which we
adopt---following Dubois-Violette--Todorov and Boyle---rather than derive from
$\Alg$; we are explicit that connecting the two algebras is an \emph{open
bridge}, not a theorem. We are equally explicit that this is a constrained
discrete model space, not a parameter-free theory: the construction rests on a
fixed complex structure and on the order of seventeen discrete algebraic choices
(Section~\ref{sec:stats}). Each component of this
claim is backed by an explicit, self-contained numerical check, collected in a
single reproducible harness with pre-registered falsification conditions
(described in Appendix~\ref{app:repro}); such a check certifies that the
finite-dimensional computation matches the stated spectrum, not that the physical
interpretation is correct.
The algebraic content (minimal left ideals, charge operators, the gauge group) is
\emph{established prior work} (Furey; Dixon; Dubois-Violette and Todorov; Boyle);
we cite it rather than claim it. Our own contribution is twofold and modest: (i)
the observation that identifying generations with the three primitive idempotents
of an $\JO$ frame---rather than with triality-permuted representations---is
\emph{not subject to} the Distler--Garibaldi obstruction \emph{at the level of
count and chirality}; and (ii) two elementary lemmas on the cross-generation Yukawa
operator that delimit precisely what the algebra does and does not fix about
flavour---in particular, that it fixes a mixing \emph{law} but selects no
hierarchy. We are explicit throughout that the \emph{continuous} flavour data (the
mass hierarchy and mixing magnitudes) is \emph{not} derived here, and we localise
the single open problem to one scalar function.
\end{abstract}

\noindent\textbf{Keywords:}\; division algebras; octonions; exceptional Jordan
algebra $\JO$; Standard Model generations; KO-dimension~6; Distler--Garibaldi
obstruction; reproducible computation.

\bigskip

\section{Scope and stance}
\label{sec:scope}

The literature on division-algebra approaches to particle physics contains a
recurring failure mode: a suggestive numerical coincidence is promoted to a
``derivation'' without an accounting of how many choices were spent to reach it.
The most prominent cautionary example is the $E_8$ embedding of
Lisi~\cite{lisi2007}, shown by Distler and Garibaldi~\cite{distler2010} to be
obstructed because a single $E_8$
cannot contain three generations of the correct chirality. We adopt the opposite
stance: we make the \emph{smallest} claim the algebra actually forces, we attach a
falsifiable check to each piece, and we keep a separate, public account of
everything we cannot yet derive (see Section~\ref{sec:not-claimed} and the
accompanying \code{METHODOLOGY\_LIMITS} and \code{PUBLIC\_CLAIMS} documents).

Concretely, this paper claims \textbf{(A)}, and explicitly disclaims \textbf{(B)}:

\begin{description}[leftmargin=2.2em,style=nextline]
\item[(A) --- claimed.] Once a complex structure is fixed, the discrete internal
  quantum numbers of one generation---charges, colour, weak isospin, hypercharge,
  anomaly freedom, and chirality---follow from $\Alg$ and its automorphism
  structure. The family count $N_{\mathrm{gen}}=3$, at the level of \emph{count and
  chirality}, and---conditionally---the absence of a fourth family, follow instead
  from the rank of an \emph{adopted} $\JO$ frame; connecting that frame to $\Alg$
  is an open bridge (Section~\ref{sec:threegen}).
\item[(B) --- disclaimed here.] The Yukawa magnitudes (fermion mass ratios), the
  mixing-angle magnitudes, the gauge couplings, the electroweak and Planck-scale
  hierarchies, and gravity are \emph{not} derived in this paper. Where the wider
  project offers candidate bridges for these, they carry open-bridge or diagnostic
  status and are out of scope for the conservative claim.
\end{description}

This division is not rhetorical. It is enforced mechanically: every statement in
Sections~\ref{sec:discrete}--\ref{sec:yukawa} corresponds to a self-contained
numerical check with a recorded status, and the claim that any item in (B) is
derived is pre-registered as a falsification condition (Appendix~\ref{app:repro}).

We state our purpose plainly at the outset, so the reader does not mistake the
scope: this paper is a \emph{delimitation}, not a derivation. Its unit of
publication is a precise, falsifiable account of what $\Alg$ and $\JO$ do and do
not fix about one generation and the family count---a boundary result that, scored
on its own diagnostics, carries \emph{no} positive evidential case against an
$O(1)$-numerology null (Section~\ref{sec:stats}). We argue in
Section~\ref{sec:conclusion} that, for a program at this stage, such a delimitation
is the correct thing to publish; we make the case here rather than only in the
conclusion.

\section{The algebra and the one-generation module}
\label{sec:algebra}

We take the internal algebra to be $\Alg = \CHO$, the tensor product of the
complex numbers, quaternions, and octonions. The construction of a single fermion
generation as a space of minimal left ideals of (the complexification of) this
algebra, and the appearance of the Standard-Model gauge group
$SU(3)_c \times SU(2)_L \times U(1)_Y$ as the subgroup of automorphisms commuting
with a fixed complex structure, are \emph{prior results} of
Furey~\cite{furey2015,furey2018} and of Dubois-Violette~\cite{duboisviolette2016},
with closely related constructions by Dixon~\cite{dixon1994},
Dubois-Violette and Todorov~\cite{todorov2018}, and Boyle~\cite{boyle2020}; for the
broader division-algebra and grand-unification background see Baez and
Huerta~\cite{baez2010}. We use
them as the foundation and do not reproduce or reclaim them. The choice of complex
structure is the step that singles out the Standard-Model gauge subgroup; we adopt
Furey's choice and do not here examine which complex structure is selected, or
why. We stress that this fix is the \emph{single most consequential} of the
discrete choices catalogued in Section~\ref{sec:stats}: it is what selects
$SU(3)_c \times SU(2)_L \times U(1)_Y$ from the automorphisms, and the conservative
claim~(A) is conditioned on it throughout. We therefore locate it explicitly in the
choice budget there rather than treat it as free, and we regard ``how constrained
is the complex structure?'' as the most important open question bearing on the
weight of~(A).

Our computational toolkit (\code{octonion\_toolkit.py}) builds the octonion
multiplication table from the Fano plane and provides the left/right
multiplication operators used throughout. All checks below run on this explicit
table; none rely on closed-form identities we have not verified numerically.

Our own contribution begins only in Sections~\ref{sec:threegen}
and~\ref{sec:yukawa}. The one-generation gauge content of
Section~\ref{sec:discrete} is inherited from the works above; we include it for
completeness and to fix notation, and we claim no novelty for it.

\section{The derived discrete structure of one generation}
\label{sec:discrete}

Each subsection states a claim, its algebraic source, and a pointer to the
self-contained numerical check that reproduces it; the checks are collected in a
reproducible harness whose conventions, and the pre-registered conditions that
would falsify each claim, are described in Appendix~\ref{app:repro} and summarised
in Table~\ref{tab:witnesses}. Each check is a reproducibility and consistency test
on an explicit finite-dimensional computation---it verifies the stated
representations, spectra, and multiplicities to machine precision and, for the
inherited results of this section, reproduces the prior-work derivations rather
than independently \emph{forcing} them---and is not, by itself, a substitute for
proof. The two genuinely novel claims of Section~\ref{sec:yukawa} are accordingly
accompanied by short proofs, not only numerical checks.

\subsection{Electric charge and colour multiplicity}

\noindent\textbf{Claim.} The one-generation electric charges are
$\{0,\tfrac13,\tfrac23,1\}$ with colour multiplicities $(1,3,3,1)$.

\noindent\textbf{Source.} The $\mathbb{C}\otimes\mathbb{O}$ number operator built
from the octonionic ladder (Furey; Dubois-Violette). The eigenvalues of the charge
operator and the dimensions of its eigenspaces are read directly off the ladder.

\noindent\textbf{Check.} \code{ladder\_charges.py} reproduces the charges and
multiplicities numerically as an \emph{output} of the algebra.
This is the standard hypercharge filter and is granted by the prior literature.

\subsection{Weak isospin and hypercharge}

\noindent\textbf{Claim.} The weak $SU(2)_L$ acts through the $\mathbb{H}$ factor,
and the Gell-Mann--Nishijima relation $Y = 2(Q - T_3)$ then fixes the full
one-generation hypercharge spectrum.

\noindent\textbf{Source.} The quaternionic factor supplies the weak doublet
structure; combined with the charges above it determines $Y$ with no free
assignment.

\noindent\textbf{Check.} \code{weak\_isospin\_hypercharge.py}
derives the hypercharges; the chiral doublet/singlet split is provided by
Section~\ref{sec:chirality}.

\subsection{Anomaly cancellation}

\noindent\textbf{Claim.} The one-generation content assembled above is free of
gauge anomalies.

\noindent\textbf{Check.} \code{physics\_map\_audit.py} freezes the quantum-number
map, verifies algebraic $Q$/$T_3$/$Y$ consistency, and checks Standard-Model
anomaly cancellation. The three-generation content map and the Yukawa spectrum are
left explicitly open.

\subsection{Chirality without fermion doubling}
\label{sec:chirality}

\noindent\textbf{Claim.} The internal space carries a chiral structure of
\emph{KO-dimension~6}---the value at which one obtains chirality without a mirror
fermion.

\noindent\textbf{Source.} Building the Clifford algebra $\Cl(0,7)$ from octonion
left-multiplications, with the real structure given by complex conjugation, the
three real-structure signs (in the Connes convention $J^2 = \epsilon$,
$JD = \epsilon' DJ$, and $J\gamma = \epsilon'' \gamma J$) come out
$(\epsilon,\epsilon',\epsilon'') = (+1,+1,-1)$, i.e.\
KO-dimension~6---the same value Connes' noncommutative-geometry Standard
Model~\cite{connes2006,chamseddine1997} requires so that the order-two axiom
removes fermion doubling.

\noindent\textbf{Check.} Verified numerically in \code{ko\_dimension\_chirality.py}.
This is the structural cure for the mirror-pair problem that otherwise sinks
generation claims (Section~\ref{sec:threegen}).

\begin{table}[t]
\centering
\footnotesize
\setlength{\tabcolsep}{4pt}
\renewcommand{\arraystretch}{1.25}
\begin{tabular}{@{}L{2.9cm}L{3.1cm}L{2.9cm}L{4.0cm}@{}}
\toprule
\textbf{Quantity} & \textbf{Result} & \textbf{Algebraic source} & \textbf{Numerical check} \\
\midrule
Electric charge / colour & $\{0,\tfrac13,\tfrac23,1\}$, mult.\ $(1,3,3,1)$ & $\mathbb{C}\otimes\mathbb{O}$ ladder & \bcode{ladder\_charges} \\
Weak isospin, hypercharge & $Y = 2(Q-T_3)$, full spectrum & $\mathbb{H}$ doublet $+$ charge & \bcode{weak\_isospin\_}\newline\bcode{hypercharge} \\
Gauge anomalies & cancel & frozen quantum-number map & \bcode{physics\_map\_audit} \\
Chirality & KO-dimension $6$, no mirror & $\Cl(0,7)$, real structure & \bcode{ko\_dimension\_}\newline\bcode{chirality} \\
Family count & $N_{\mathrm{gen}}=3$ (count $+$ chirality) & rank of $\JO$, inner $S_3<F_4$ & \bcode{jordan\_eigenvalue\_}\newline\bcode{generations},\newline\bcode{three\_generations\_}\newline\bcode{frame} \\
No fourth family & excluded (conditional) & $H_n(\mathbb{O})$ Jordan only for $n\le 3$ & conditional theorem \\
\bottomrule
\end{tabular}
\caption{The discrete internal structure of one generation. Each row is an
\emph{output} of the algebra and its automorphisms; the first five are verified by
a self-contained numerical check (Appendix~\ref{app:repro}), while the exclusion
of a fourth family is a conditional theorem. None requires a per-observable
continuous input.}
\label{tab:witnesses}
\end{table}

\section{The family count: three generations, count and chirality}
\label{sec:threegen}

The headline integer of the framework is $N_{\mathrm{gen}}=3$. We are careful to
claim it at exactly the level it is established---\emph{count and chirality}---and
no further.

\paragraph{Two algebras, and the bridge between them.}
Sections~\ref{sec:algebra}--\ref{sec:discrete} build one generation
inside $\Alg = \CHO$, whose relevant module is
$\mathbb{C}\otimes\mathbb{O} = \mathbb{C}^8$ ($16$ real dimensions). The
generation \emph{count} of this section, by contrast, lives in the exceptional
Jordan algebra $\JO$ ($27$ real dimensions), which is \emph{not} a subalgebra or a
tensor factor of $\Alg$---indeed $27 \nmid 64$. We do not derive $\JO$ from
$\Alg$; we adopt it as an additional algebraic structure, following
Dubois-Violette--Todorov and Boyle, motivated by the observation below that the
tangent space at each of its primitive idempotents is a real $\Spin(9)$ spinor
$\Delta_9 \cong \mathbb{C}\otimes\mathbb{O}$---the same $16$-dimensional
one-generation module. The integer ``three'' is therefore the \emph{rank} of
$\JO$; reading it as the Standard-Model family count, and identifying $\JO$ with
the generation structure of $\Alg$, is an \emph{open bridge}, not a theorem. We
flag this prominently because the claim is otherwise easily over-read: the count
follows from positing a rank-three object, and what $\Alg$ contributes is the
one-generation module that sits at each idempotent.

\subsection{Why the naive route fails}

The intuitive identification ``three triality representations $=$ three
generations'' is \emph{obstructed}, and we say so first. Our own stress test
(\code{three\_generations\_nogo\_audit.py}) reproduces the Distler--Garibaldi
failure mode: the three 8-dimensional $\Spin(8)$ representations are a vector
$8_v$ and a chirality-mirror pair of spinors $8_s, 8_c$. Identifying generations
with them would (i) require the outer triality to map a vector to a spinor, and
(ii) hand each generation a mirror partner. This is the same obstruction that
defeats the $E_8$ proposal~\cite{distler2010}, and it downgrades the naive bridge
to conjecture.

\subsection{The idempotent-frame resolution}

The resolution is to identify the three generations \emph{not} with
triality-permuted representations but with the three \emph{primitive idempotents}
of a maximal frame of the exceptional Jordan algebra $\JO$---three identical
points of the Cayley plane $\OP = F_4/\Spin(9)$. Two independent results support
this.

\paragraph{Spectral count.}
\code{jordan\_eigenvalue\_generations.py} verifies that the Freudenthal
characteristic polynomial of a Hermitian $\JO$ element is a cubic with three real
roots and three orthogonal primitive idempotents resolving the identity (checked
on 4000 random samples). ``Three'' is the \emph{rank} of the algebra---the degree
of the cubic norm---and is therefore immune to the representation-theoretic mirror
obstruction entirely.

\paragraph{Frame equivalence and chirality.}
\code{three\_generations\_frame.py} builds $\ffour = \Der(\JO)$ (dimension 52) and
the idempotent stabiliser $\mathfrak{spin}(9)$ (dimension 36, semisimple) from the
octonion table, and shows the three idempotents are permuted by the \emph{inner}
frame $S_3 < F_4$. Because $F_4$ is connected with no outer automorphism, this
$S_3$ cannot carry a representation to an inequivalent one, so Obstruction~(i)
cannot even be posed. The 16-dimensional tangent space at each idempotent is a
single \emph{real} $\Spin(9)$ spinor $\Delta_9$ (the octonionic $\Cl(9)$ module,
commutant dimension 1)---identical to the KO-dimension-6 module of
Section~\ref{sec:chirality}---so three identical self-conjugate copies share one
chirality with no mirror partner, and Obstruction~(ii) cannot arise either.

Because this last point carries real weight---it is both the chirality cure of
Section~\ref{sec:chirality} and the evasion of Obstruction~(ii)---we give the
argument explicitly rather than leave it to the numerical check. $\Spin(9)$ is the spin
group of an \emph{odd}-dimensional space, so it has a \emph{single} irreducible
spinor representation $\Delta_9$: there is no $\Delta_9^{L}/\Delta_9^{R}$ chirality
splitting of the kind that exists in even dimension. Moreover $\Delta_9$ is of
\emph{real} type---its commutant is $\mathbb{R}$ (dimension~1, as the numerical
check verifies)---so it admits a $\Spin(9)$-invariant real structure and is
isomorphic to its own conjugate, $\overline{\Delta_9} \cong \Delta_9$. A mirror
partner would have to be an \emph{inequivalent} conjugate (opposite-chirality)
module; for a single self-conjugate real representation no such module exists.
Hence three identical copies of $\Delta_9$ cannot furnish mirror fermions,
independently of any content map.

The $8_v/8_s/8_c$ triple that the no-go attacks is shown (Part~D of the same
module) to be the \emph{off-diagonal} Peirce decomposition, a genuinely different
object that does mix chirality---confirming the obstruction is real for the naive
route and absent for the idempotent route.

\paragraph{Relation to Boyle and Dubois-Violette--Todorov.} The closest prior work
places one generation in (the tangent space of) the complex exceptional Jordan
algebra and relates the \emph{three} generations to $SO(8)$ triality:
Boyle~\cite{boyle2020} states exactly this, and Dubois-Violette--Todorov
\cite{duboisviolette2016,todorov2018} locate the three families in the three
off-diagonal octonionic directions of $\JO$. By the analysis above, those are
precisely the objects ($8_v, 8_s, 8_c$; the off-diagonal Peirce slots) on which
the Distler--Garibaldi obstruction bites. Our contribution is narrow and specific:
the \emph{same count} is available \emph{without} triality, from the three
\emph{diagonal} primitive idempotents of a single real $\JO$ frame, permuted by an
\emph{inner} $S_3 \subset F_4$. We do not claim a different or better physical
generation structure than these works---only a different \emph{route to the count}
that is not subject to the obstruction \emph{at the level of count and chirality}.
We are careful not to overstate this: the Distler--Garibaldi obstruction is
ultimately about realising three generations \emph{with the correct chiral gauge
action}, and that fermion-content map is exactly what we leave open below. Our
route is therefore better described as \emph{not yet engaging} the obstruction on
the content map than as \emph{defeating} it there; what we do establish is that the
count and the shared chirality are obstruction-free.
The price is real, and we state it: because the three
idempotents are $F_4$-equivalent, this structure carries \emph{no} flavour
information beyond the integer three (all flavour must come from the
symmetry-breaking of Section~\ref{sec:yukawa})---arguably weaker, in that respect,
than triality routes in which $8_v, 8_s, 8_c$ are at least distinct objects.

\subsection{No fourth generation}

This is the most conditional of our claims, and we flag it as such. Within the
present construction the family count is the \emph{rank} of the Jordan algebra,
and that rank is capped at three for a structural reason: the Hermitian
$n \times n$ octonionic matrices $H_n(\mathbb{O})$ satisfy the Jordan axioms only
for $n \le 3$ (Jordan--von Neumann--Wigner~\cite{jvnw1934}); there is no
exceptional Jordan algebra $\mathfrak{J}_4(\mathbb{O})$, and the naive extension
passes to the non-alternative sedenions. A fourth family is therefore excluded
\emph{relative to} the identification of generations with a primitive-idempotent
frame of $\mathfrak{J}_3(\mathbb{O})$. We state it as a conditional
exclusion, not an unconditional no-go: it would not bind if generations were
realised by some construction other than such a frame. Because the cap follows
from the \emph{same} rank-three identification that produces the count, it carries
essentially no \emph{independent} predictive content, and we do not present it as a
separate success of the framework.

\subsection{What this section does \emph{not} claim}

We do not claim the \emph{fermion-content map}: the assignment of the $16$ real
dimensions of $\Delta_9$ to the $16$ Weyl fields of one Standard-Model generation,
\emph{with their correct chiral gauge action}, is not established here. Until that
map exists, the result of this section is a counting statement about the three
primitive idempotents of a $\JO$ frame---together with the verified fact that their
tangent spaces are three identical real $\Spin(9)$ spinors---rather than a
statement about three Standard-Model generations. A flavour-diagonal Jordan
element still returns only its seeded diagonal (a genuine negative). The count and
the chirality \emph{type} are obstruction-free; the content map and the
\emph{magnitudes} that distinguish the generations are the subject of
Section~\ref{sec:yukawa} and remain open.

\section{The cross-generation Yukawa operator: two elementary lemmas}
\label{sec:yukawa}

Because the three generations are \emph{identical} points of $\OP$, nothing in
Section~\ref{sec:threegen} distinguishes them---the mass hierarchy and mixing live
entirely in the Yukawa operator. Here we record two elementary lemmas that delimit
what the algebra fixes about that operator. Both are immediate from the diagonal
action on the Peirce slots (the proofs are two lines each); we state them
carefully because their \emph{net} content---that the algebra fixes a mixing law
but selects no hierarchy---is the boundary of what $\JO$ contributes to
flavour.

\subsection{Setup}

A single algebra-internal Dirac operator on the one-generation module
$\mathbb{C}\otimes\mathbb{O} = \mathbb{C}^8$ is isospectral and carries \emph{no}
forced mass ratio (the elimination argument of \code{spectral\_action.py}); a
generation spectrum must therefore come from a structure that carries the three
families, which in the present setting is the $\JO$ frame of
Section~\ref{sec:threegen}---with the three generations as its primitive
idempotents, subject to the open bridge flagged there. The natural algebra-internal Yukawa operators on $\JO$ are the Jordan
\emph{left multiplication} $L_X : Y \mapsto X \circ Y$ and the canonical Jordan
\emph{quadratic representation} $U_X : Y \mapsto 2X\circ(X\circ Y) - (X\circ X)\circ Y$.
We build both from the explicit $\JO$ structure tensor.

\subsection{Lemma 1 --- the averaging law (additive mixing)}

\begin{lemma}[Averaging law]
\label{thm:avg}
For $X = \mathrm{diag}(a,b,c)$, the spectrum of $L_X$ is exactly $\{a,b,c\}$
together with $\bigl\{\tfrac{a+b}{2}, \tfrac{b+c}{2}, \tfrac{c+a}{2}\bigr\}$, the
latter each with multiplicity $8$. Every inter-generation level is the
\emph{arithmetic mean} of two generation levels.
\end{lemma}

\begin{proof}[Proof sketch]
In the diagonal frame, $\JO$ decomposes as the three real diagonal entries
together with the three off-diagonal octonionic slots $z_1, z_2, z_3$ (each
$\cong \mathbb{O} \cong \mathbb{R}^8$), giving $3 + 24 = 27$. For
$X = \mathrm{diag}(a,b,c)$ the Jordan product $X \circ Y = \tfrac12(XY + YX)$ acts
on a diagonal entry $y_{ii}$ as multiplication by $x_i$, and on a slot entry
$y_{ij}$ ($i \neq j$) as
\[
  (X \circ Y)_{ij} = \tfrac12\bigl(x_i\, y_{ij} + y_{ij}\, x_j\bigr)
                   = \tfrac12 (x_i + x_j)\, y_{ij},
\]
since the $x_i$ are real and commute with octonions. Hence $L_X$ is block-diagonal
in this decomposition, with eigenvalues $a,b,c$ on the diagonal entries and
$\tfrac{a+b}{2}, \tfrac{b+c}{2}, \tfrac{c+a}{2}$ on the three $8$-dimensional
slots. The averaging law and the multiplicities $[1,1,1,8,8,8]$ follow.
\end{proof}

\noindent\textbf{Check.} \code{spectral\_action\_432.py} verifies this over 200
random seeds with residual $\sim 10^{-7}$ and the multiplicity pattern
$[1,1,1,8,8,8]$. This is three parameter-free relations among the 27 eigenvalues,
holding for \emph{all} $X$.

\noindent\textbf{Consequence.} The single existing triality-breaking spurion
$\epso^2 = \pi/432$ (used here only as a symmetry-breaking \emph{scale}; we do
\emph{not} claim its value is geometrically forced---see
Section~\ref{sec:not-claimed}) breaks the inner $S_3$ and reduces the generation
freedom from
three eigenvalues to \emph{one} scale. But because $L_X$ mixing is \emph{additive}
(arithmetic means), a one-knob $\epso$ ladder cannot reproduce a steep
multiplicative hierarchy: it misses the charged-lepton spectrum by $\approx 1.4$
decades on the lightest state. The open problem is thereby \emph{localised to one
scalar seed function}---the profile of the three diagonal eigenvalues---with the
mixing law itself derived.

\subsection{Lemma 2 --- multiplicative mixing from the quadratic operator}

\begin{lemma}[Multiplicative mixing]
\label{thm:mult}
For $X = \mathrm{diag}(a,b,c)$, the spectrum of the canonical quadratic
representation $U_X$ is exactly $\{a^2,b^2,c^2\}$ together with
$\{ab, bc, ca\}$, the latter each with multiplicity $8$. Every inter-generation
level is the \emph{geometric mean} (product) of two generation levels.
\end{lemma}

\begin{proof}[Proof sketch]
Write $s = \tfrac12(x_i + x_j)$ for the scalar by which $X \circ (\cdot)$ acts on
the slot $y_{ij}$ (proof of Lemma~\ref{thm:avg}), and note that
$X \circ X = \mathrm{diag}(a^2,b^2,c^2)$ acts there as $t = \tfrac12(x_i^2 +
x_j^2)$. The quadratic representation $U_X Y = 2\,X\circ(X\circ Y) - (X\circ
X)\circ Y$ then acts on $y_{ij}$ as
\[
  2 s^2 - t = \tfrac12 (x_i + x_j)^2 - \tfrac12 (x_i^2 + x_j^2) = x_i x_j,
\]
and as $x_i^2$ on the diagonal entries. The product (geometric-mean) law and the
multiplicities $[1,1,1,8,8,8]$ follow.
\end{proof}

\noindent\textbf{Check.} \code{spurion\_perturbation.py} verifies this over 200
random seeds with residual \emph{exactly} $0$ and multiplicity $[1,1,1,8,8,8]$.

\noindent\textbf{Significance.} Under $U_X$, log-mass is \emph{additive} in the
generation exponents. Multiplicative mixing is the structural prerequisite for any
power-law hierarchy; switching from the linear Yukawa $L_X$ to the canonical Jordan
quadratic $U_X$ is therefore the algebraic step that \emph{permits} a power-law
ladder at all. We are candid that the \emph{motivation} for preferring $U_X$ is
phenomenological---one wants multiplicativity because the observed hierarchy is
roughly power-law; the lemma itself, however, is a fixed property of the operator,
independent of any data. Neither lemma \emph{selects} a hierarchy: the off-diagonal
levels are means or products of $(a,b,c)$, so all hierarchy must be seeded in the
diagonal profile, which the algebra leaves unexplained.

\subsection{A companion structural fact, and the boundary}

One further observation sharpens---but does \emph{not} close---the magnitude
problem.

\begin{itemize}[leftmargin=1.4em]
\item \textbf{Rank-one bottleneck.} The project's single spurion is rank one,
  $T_{\mathrm{break}} = \theta\,|\tau\rangle\langle\tau|$. A rank-one perturbation
  of a degenerate level lifts exactly one eigenvalue at first order (verified
  numerically in \code{spurion\_perturbation.py}). Hence a three-tier hierarchy
  cannot arise from one first-order insertion; the tiers must appear at cumulative
  orders $\epso^1, \epso^2, \dots$
\end{itemize}

We deliberately exclude from this paper any numerical pattern in the diagonal
exponents---for example the $\epso = \sqrt{\pi/432}$ ladder explored in the wider
project---because such patterns are post-hoc targets over small, self-chosen
windows and would invite the very look-elsewhere critique the conservative claim is
designed to avoid. We refer them to the companion materials and do not display them
here.

The boundary of this paper is exactly here: Section~\ref{sec:yukawa} fixes
the \emph{mixing law} and the \emph{operator class}, and localises the residue to
one diagonal seed profile. The \emph{dynamical selection} of that profile is not
supplied by the algebra and is the single genuinely open research item. We state
it as such rather than fitting it.

\section{Statistical assessment}
\label{sec:stats}

A discrete-structure claim must be defended against the charge that integers are
cheap. We therefore report the project's adversarial diagnostics rather than
suppress them.

\begin{itemize}[leftmargin=1.4em]
\item \textbf{Parameter count (MDL).} The framework carries roughly 17 discrete
  structural choices plus one continuous input (the Planck mass), not zero
  (\code{model\_complexity.py}). The data-compression ratio is marginal today
  ($R \approx 1.19$).
\item \textbf{Goodness-of-fit on independent observables.} Once dependent rows are
  removed and a stated theory floor is applied, the independent set is
  statistically consistent (reduced $\chi^2 \approx 0.92$), with the
  first-generation electron mass a visible outlier
  (\code{independent\_observables.py}, \code{covariance\_gof.py}). The 22 nominal
  rows collapse to $N_{\mathrm{eff}} \approx 10$ effective observables under the
  shared-$\epso$ common mode---we do not count them as 22 independent successes.
\item \textbf{Model-comparison Bayes factor.} Crediting only the
  numerically-closed discrete results (the conservative floor of this paper), the
  Bayes factor against an $O(1)$-numerology null is approximately
  $\ln B \approx -3.2$; it rises to $\approx +5.6$ only if the geometric origin of
  $\pi/432$ is granted (\code{bayesian\_evidence.py}, \code{scoreboard.py}). The
  verdict
  therefore hinges on a single named seam, which we keep on the disclaimed side
  (B). The conservative paper does \emph{not} rely on $\ln B > 0$.
\item \textbf{Look-elsewhere on scale relations.} The power-of-three scale hits
  ($M_W$, $M_R$, $\Lambda$) are \emph{cheap}: a simple prefactor times an integer
  exponent already covers $\approx 93\%$ of one exponent window
  (\code{scale\_look\_elsewhere.py}). These are explicitly \emph{not} part of claim
  (A).
\end{itemize}

The conservative claim survives all of these because it rests on \emph{discrete},
structure-level results (charges, multiplicities, anomaly cancellation,
KO-dimension, the rank-three count) whose evidential weight does not depend on the
continuous-parameter accounting that the diagnostics (correctly) discount.

\section{What this paper does not claim}
\label{sec:not-claimed}

For the avoidance of doubt, the following are \emph{not} claimed here and are
pre-registered as falsification conditions (Appendix~\ref{app:repro}):

\begin{enumerate}[leftmargin=1.8em]
\item The fermion mass ratios / Yukawa eigenvalue magnitudes are not derived.
\item The mixing-angle magnitudes are not derived; the integer multiplicities
  $\sqrt{7}, \tfrac12, 3, 4$ and $4/7$ in the wider project are open bridges, not
  closed.
\item $\epso^2 = \pi/432$ is not claimed as geometrically forced; it is an open
  bridge.
\item The gauge couplings and the electroweak / Planck-scale hierarchies are not
  derived; the relevant rows carry continuum/RG-matching gaps obtained partly by
  inverse running.
\item Gravity is out of scope: the internal metric brick yields no canonical 4D
  Lorentzian reduction or dynamics (\code{gravity\_gate\_audit.py}).
\end{enumerate}

\section{Conclusion}
\label{sec:conclusion}

The defensible \emph{content} of the $\CHO$ program is sharp and small: the
\emph{discrete internal anatomy of one Standard-Model generation}---charges,
colour, weak isospin, hypercharge, anomaly freedom, and chirality---is determined,
once a complex structure is fixed, by the algebra and its automorphisms; and the
family count $N_{\mathrm{gen}}=3$, at the level of count and chirality, is the
rank of an $\JO$ frame whose primitive idempotents are permuted by an inner $S_3$
---a route that is not subject to the
Distler--Garibaldi no-go that defeats representation-counting bridges.
We are explicit that this last step rests on adopting $\JO$ (an open
$\Alg \to \JO$ bridge and an unestablished fermion-content map) and that the
$F_4$-equivalent idempotents carry no flavour information beyond the integer
three.

Our new contribution to the flavour question is delimitation rather than
solution: the cross-generation Yukawa obeys an \emph{averaging law} under linear Jordan
multiplication (Lemma~\ref{thm:avg}) and a \emph{multiplicative-mixing law} under
the canonical quadratic operator (Lemma~\ref{thm:mult}), and the rank-one spurion
forces any hierarchy to be cumulative. These results reduce the open flavour problem from ``derive the entire
Yukawa spectrum'' to ``supply a dynamical principle selecting one diagonal seed
profile.'' That residual is stated plainly as the program's lone high-risk open
item---not papered over with a fit.

We stress, finally, what this paper does \emph{not} establish: a \emph{positive}
evidential case. Scored on its own diagnostics (Section~\ref{sec:stats}), the
conservative results do not beat an $O(1)$-numerology null ($\ln B \approx -3.2$)
unless the
geometric origin of $\pi/432$ is granted---which we keep on the disclaimed side
(B). The content here is therefore offered as a precise, falsifiable
\emph{delimitation} of what $\CHO$ and $\JO$ do and do not fix, not as positive
evidence for the framework. We regard that delimitation as the correct unit of
publication: every component is individually falsifiable, the inherited results
are attributed rather than reclaimed, and the open bridges are named rather than
papered over.

\section*{Code and data availability}
\addcontentsline{toc}{section}{Code and data availability}

All claims in this paper are checked by a single reproducible harness that
accompanies this deposit and is openly available in the project repository,
\url{https://github.com/richorama/cho}. The numerical-check modules of
Sections~\ref{sec:discrete}--\ref{sec:yukawa} are:

\begin{quote}\ttfamily\small\raggedright
ladder\_charges.py, weak\_isospin\_hypercharge.py, physics\_map\_audit.py,
ko\_dimension\_chirality.py, jordan\_eigenvalue\_generations.py,
three\_generations\_frame.py, spectral\_action\_432.py, spurion\_perturbation.py
\end{quote}

\noindent together with the statistical diagnostics of Section~\ref{sec:stats}.
They reside in the \code{compute/} directory of that repository. The full
suite, the per-claim contracts, and the test harness run with:

\begin{quote}\ttfamily\small
python3 compute/audit.py \quad\textrm{\# run every artifact}\\
python3 compute/audit.py <name> \quad\textrm{\# run one}\\
python3 compute/audit\_contract.py \quad\textrm{\# validate claim contracts}\\
python3 -m pytest -q \quad\textrm{\# one test per audit artifact}
\end{quote}

The numerical inputs use NumPy only. The status of every claim is pinned in the
companion ledger (\code{DERIVATION\_LEDGER.md}); the public-claim policy is in
\code{PUBLIC\_CLAIMS.md}; caveats are in \code{METHODOLOGY\_LIMITS.md}.

\paragraph{License.}
This work is released under the Creative Commons Attribution 4.0 International
(CC~BY~4.0) license.

\appendix

\section{Reproducibility apparatus}
\label{app:repro}

The numerical checks cited in the body are organised by a small, project-internal
bookkeeping layer. We collect its conventions here so that the main text can refer
to them in plain language; none of this vocabulary is required to follow the
physics.

\paragraph{Vocabulary.}
A \emph{check} (a \emph{witness}, in the repository) is the executable script that
reproduces a stated claim on the explicit octonion multiplication table. The
\emph{ledger} (\code{DERIVATION\_LEDGER.md}) records the status of each claim. A
\emph{kill condition} is a pre-registered statement whose truth would falsify a
claim, and a \emph{contract} ties each check to its ledger status---so that
asserting any disclaimed item (B) is derived would itself fail the harness. A
check returning \textsc{pass} certifies only that the finite-dimensional
computation matches the stated representations, spectra, and multiplicities to
machine precision; it is not a proof and says nothing about the physical
interpretation.

\paragraph{Internal labels.}
``Lever A/B/C/D'' are repository labels for four one-generation results: the
spectral generation count~(A), the KO-dimension-6 chirality~(B), the charge
ladder~(C), and weak isospin and hypercharge~(D). They are retained for
traceability only and carry no content beyond the results they name.

\paragraph{Per-claim checks.}
The body pointers map to modules as follows: electric charge and colour,
\code{ladder\_charges.py}; weak isospin and hypercharge,
\code{weak\_isospin\_hypercharge.py}; anomaly cancellation,
\code{physics\_map\_audit.py}; KO-dimension-6 chirality,
\code{ko\_dimension\_chirality.py}; the spectral generation count,
\code{jordan\_eigenvalue\_generations.py}; frame equivalence and chirality,
\code{three\_generations\_frame.py}; the averaging law (Lemma~\ref{thm:avg}),
\code{spectral\_action\_432.py}; and multiplicative mixing (Lemma~\ref{thm:mult})
together with the rank-one bottleneck, \code{spurion\_perturbation.py}. The
exclusion of a fourth family is the conditional ledger entry~G2, with no separate
module. Run instructions are given in the \emph{Code and data availability}
section above.

\begin{thebibliography}{11}

\bibitem{lisi2007}
A.~G.~Lisi,
\emph{An Exceptionally Simple Theory of Everything},
\texttt{arXiv:0711.0770} (2007).

\bibitem{furey2015}
C.~Furey,
\emph{Standard model physics from an algebra?},
Ph.D.\ thesis, University of Waterloo (2015); \texttt{arXiv:1611.09182} (2016).

\bibitem{furey2018}
C.~Furey,
\emph{Three generations, two unbroken gauge symmetries, and one
eight-dimensional algebra},
Phys.\ Lett.\ B \textbf{785} (2018) 84--89, \texttt{arXiv:1910.08395}.

\bibitem{dixon1994}
G.~M.~Dixon,
\emph{Division Algebras: Octonions, Quaternions, Complex Numbers and the
Algebraic Design of Physics},
Kluwer Academic Publishers (1994).

\bibitem{duboisviolette2016}
M.~Dubois-Violette,
\emph{Exceptional quantum geometry and particle physics},
Nucl.\ Phys.\ B \textbf{912} (2016) 426--449, \texttt{arXiv:1604.01247}.

\bibitem{todorov2018}
M.~Dubois-Violette and I.~Todorov,
\emph{Exceptional quantum geometry and particle physics II},
Nucl.\ Phys.\ B \textbf{938} (2019) 751--761, \texttt{arXiv:1808.08110}.

\bibitem{baez2010}
J.~C.~Baez and J.~Huerta,
\emph{The algebra of grand unified theories},
Bull.\ Amer.\ Math.\ Soc.\ \textbf{47} (2010) 483--552, \texttt{arXiv:0904.1556}.

\bibitem{boyle2020}
L.~Boyle,
\emph{The Standard Model, the Exceptional Jordan Algebra, and Triality},
\texttt{arXiv:2006.16265} (2020).

\bibitem{jvnw1934}
P.~Jordan, J.~von Neumann, and E.~Wigner,
\emph{On an algebraic generalization of the quantum mechanical formalism},
Ann.\ Math.\ \textbf{35} (1934) 29--64.

\bibitem{connes2006}
A.~Connes,
\emph{Noncommutative geometry and the standard model with neutrino mixing},
JHEP \textbf{0611} (2006) 081, \texttt{arXiv:hep-th/0608226}.

\bibitem{chamseddine1997}
A.~H.~Chamseddine and A.~Connes,
\emph{The Spectral Action Principle},
Commun.\ Math.\ Phys.\ \textbf{186} (1997) 731--750, \texttt{arXiv:hep-th/9606001}.

\bibitem{distler2010}
J.~Distler and S.~Garibaldi,
\emph{There is no ``Theory of Everything'' inside $E_8$},
Commun.\ Math.\ Phys.\ \textbf{298} (2010) 419--436, \texttt{arXiv:0905.2658}.

\end{thebibliography}

\end{document}
